% chapters/ch8_exp_scan70.tex --- 4.4 受信条件悪化時の頑健性検証実験
\section{受信条件悪化時の頑健性検証実験}
\label{sec:exp_scan70}

\subsection{実験の目的}

スキャンデューティ比を70\%に低下させた条件で，提案手法の頑健性を検証する．

\subsection{実験結果}

\tabref{tab:d3_summary}に結果を示す．固定\SI{500}{\milli\second}は$P_{\mathrm{out}}=0.285$と悪化したが，動的制御は0.089と改善し，固定\SI{100}{\milli\second}（0.065）と同様に制約$\epsilon=0.1$を満たした．

\begin{table}[tb]
  \centering
  \caption{scan70での結果（n=3）}
  \begin{tabular}{lrrrr}
    \toprule
    条件 & $P_{\mathrm{out}}(1\mathrm{s})$ & P[mW] & $q_{\mathrm{event}}$ \\
    \midrule
    100ms & 0.065 & 209.9 & 281.4 \\
    500ms & 0.285 & 189.5 & 251.6 \\
    policy & 0.089 & 202.1 & 270.2 \\
    \bottomrule
  \end{tabular}
\end{table}

\clearpage
